% SOURCE_STATEMENT: "Does the six-vertex model on a torus undergo a phase transition at which the emergent torus geometry collapses, with critical modular parameter $\\tau$ satisfying $\\operatorname{Im}(\\tau)\\neq0$?"
% Preparation and final self-audit: CONTRIBUTING.md prompts C and D; source/audit artifacts retained in the local physics_statements_2026_10_11 batch.
% !TeX program = lualatex
\documentclass[11pt]{article}
\usepackage[margin=25mm]{geometry}
\usepackage{fontspec}
\setmainfont{Latin Modern Roman}
\usepackage{amsmath,amssymb,mathtools}
\usepackage{unicode-math}
\setmathfont{Latin Modern Math}
\usepackage{xurl,hyperref}
\hypersetup{colorlinks=true,urlcolor=blue}
\setcounter{secnumdepth}{0}
\setlength{\parindent}{0pt}
\setlength{\parskip}{5pt}
\setlength{\emergencystretch}{3em}
\newcommand{\sref}[2]{\hyperlink{source:#1:#2}{[S#2]}}
\newcommand{\cataloguescope}[1]{\paragraph{Recorded target.}#1}
\newcommand{\R}{\mathbb R}
\newcommand{\C}{\mathbb C}
\newcommand{\N}{\mathbb N}
\newcommand{\Z}{\mathbb Z}
\newcommand{\Q}{\mathbb Q}
\newcommand{\F}{\mathbb F}
\newcommand{\eps}{\varepsilon}
\begin{document}
\section{Geometric Collapse in the Toroidal Six-Vertex Model}\label{statement}

\subsection{Definitions and mathematical statement}

\paragraph{Shared notation.}$\mathbb N=\{1,2,\ldots\}$; the standard number fields have their usual meanings. All additional conventions are stated locally.


\paragraph{Definitions and conventions.} In the six-vertex model, arrows on edges of a square lattice obey the ice rule: two arrows enter and two leave each vertex. Periodic boundary conditions place the lattice on a torus. The report discusses a continuum geometry emerging from finite-size spectral data, described by a complex modular parameter $\tau=\tau_0+i\tau_1$ with $\tau_1>0$. Geometric collapse refers to degeneration of that emergent geometry at a thermal phase transition, rather than removal of the combinatorial periodic lattice.

\paragraph{Source scope and qualifications.} The comparison comes from a dimer model where the torus angle collapses as the imaginary part tends to zero. The six-vertex conjecture deliberately predicts different behavior with nonzero critical imaginary part; it cannot be replaced by the usual cusp degeneration. The exact geometric-collapse observable remains part of the source construction.

\paragraph{Catalogue mathematical target.}
\paragraph{Catalogue wording retained for traceability.}
Does the six-vertex model on a torus undergo a phase transition at which the emergent torus geometry collapses, with critical modular parameter $\tau$ satisfying $\operatorname{Im}(\tau)\neq0$?

\paragraph{Saved literature context.} The supplied catalogue's assessment dated 2026-08-21 records: Rigorous torus free-energy and finite-size results now exist for the six-vertex model, but the specific nonzero-critical-parameter collapse conjecture is not established. The assessment is retained as dated provenance; it is not a new comprehensive literature or open-status certification.
\subsection{Short English statement}
Test the proposed six-vertex phase transition in which emergent torus geometry collapses at a critical modulus with nonzero imaginary part.
\subsection{Sources}
\begin{itemize}
\item[\hypertarget{source:statement:1}{S1}] ULAM.AI and the UnsolvedMath Contributors, UnsolvedMath, catalogue entry. Catalogue curation under CC BY 4.0; definitions and English exposition added locally.
\url{https://huggingface.co/datasets/ulamai/UnsolvedMath}.
\item[\hypertarget{source:statement:2}{S2}] Noncommutative Geometry (2009), Oberwolfach Reports, DOI 10.4171/owr/2009/41, pp. 2271--2334; official report source and mathematical context.
\url{https://doi.org/10.4171/owr/2009/41}.
\item[\hypertarget{source:statement:3}{S3}] Official EMS Press PDF of the same Oberwolfach report; PDF page 44, printed page 2314.
\url{https://ems.press/content/serial-article-files/46243}.
\item[\hypertarget{source:statement:4}{S4}] six-vertex free energy; reference in the saved dated catalogue literature review.
\url{https://arxiv.org/abs/2012.11675}.
\item[\hypertarget{source:statement:5}{S5}] toroidal finite-size partition function; reference in the saved dated catalogue literature review.
\url{https://doi.org/10.1103/PhysRevE.53.3225}.
\end{itemize}
\label{end:statement}
\end{document}
